\section{Metodologija}
\label{sec:metoda}

\subsection{Eksperimentalni postav}

Mjerenje je provedeno na stereovizijskom sustavu sastavljenom od dviju kamera
usmjerenih prema površini ispitnog uzorka pod kutom (Slika \ref{fig:overlay}).
Uzorak je vlačna epruveta s reduciranim mjernim presjekom, čija je mjerna površina
zakrivljena (vidljivo iz rekonstruiranog oblika na Slici \ref{fig:dispmag}); mjerna
zona prethodno je pripremljena nanošenjem nasumične crno-bijele teksture nužne za
korelaciju. Tijekom pokusa obje su kamere snimile po 1188 vremenski usklađenih slika,
od kojih prva odgovara nedeformiranom (referentnom) stanju, a ostale rastućim razinama
deformacije. U ovom je radu analizirano prvih 300 deformacijskih stanja.

\subsection{Kalibracija i geometrijski model}

Stereovizijski sustav kalibriran je unaprijed, čime su određene projekcijske matrice
obiju kamera $\mathsf{P}_1$ i $\mathsf{P}_2$ te trodimenzijska mreža konačnih
elemenata $\mathsf{M}$ koja opisuje referentni oblik mjerne površine (252 čvora, 438
trokutastih elemenata). Ispravnost kalibracije provjerena je projiciranjem čvorova 3D
mreže u obje kamere prema \eqref{eq:projection}; kao što pokazuje Slika
\ref{fig:overlay}, projicirana mreža precizno prekriva teksturiranu mjernu zonu u obje
slike, čime je potvrđena geometrijska konzistentnost sustava.

\begin{figure}[H]
  \centering
  \includegraphics[width=0.49\textwidth]{overlay_cam1.png}
  \includegraphics[width=0.49\textwidth]{overlay_cam2.png}
  \caption{Provjera kalibracije: 3D mreža projicirana matricama $\mathsf{P}_1$ (kamera
  1, lijevo) i $\mathsf{P}_2$ (kamera 2, desno) na referentnu sliku. Mreža prekriva
  teksturiranu mjernu zonu epruvete.}
  \label{fig:overlay}
\end{figure}

\subsection{Postupak mjerenja}

Primijenjeni postupak stereo korelacije provodi se u sljedećim koracima:
\begin{enumerate}
  \item \textbf{2D korelacija po kameri.} Za svaku kameru zasebno provodi se globalna
  korelacija konačnih elemenata između referentne i svake deformirane slike, čime se
  dobiva polje ravninskih (2D) pomaka čvorova mreže $\mathbf{U}_1$ i $\mathbf{U}_2$.
  \item \textbf{Korespondencija čvorova.} Budući da su obje slikovne mreže projekcije
  iste 3D mreže $\mathsf{M}$, korespondencija čvorova između kamera je zajamčena po
  konstrukciji -- čvor $i$ u kameri 1 odgovara čvoru $i$ u kameri 2.
  \item \textbf{Stereo triangulacija.} Za svaki frame i svaki čvor rekonstruira se 3D
  položaj triangulacijom \eqref{eq:dlt} iz njegovih deformiranih položaja u objema
  kamerama. Referentni 3D oblik dobiva se triangulacijom nedeformiranih položaja.
  \item \textbf{Polje 3D pomaka.} Trodimenzijski pomak čvora jednak je razlici
  trianguliranog deformiranog i referentnog položaja,
  $\mathbf{U}^{\text{3D}} = \mathbf{X}^{\text{def}} - \mathbf{X}^{\text{ref}}$.
  \item \textbf{Deformacije.} Iz polja 3D pomaka računaju se površinske deformacije
  \eqref{eq:green-lagrange} po elementu.
\end{enumerate}

Ovakav pristup -- 2D korelacija po kameri praćena triangulacijom -- omogućuje
izravno ponovno korištenje postojećeg, provjerenog algoritma 2D korelacije, uz
minimalnu dodatnu programsku logiku za stereo rekonstrukciju.
